\subsection{CAUCHY FILTERS}

Once a set X has been endowed with a uniform structure we can define what is meant by a ``small'' subset of X (relative to this structure): a ``small'' subset of X is one in which all the points are ``very close'' to each other. Precisely:

DEFINITION 1. If X is a uniform space and if V is an entourage of X, a subset A of X is said to be V-small if every pair of points of A are V-close (in other words, if A × A ⊂ V).

PROPOSITION 1. In a uniform space X, if two sets A and B are V-small and intersect, then their union A ∪ B is V-small.

Let \(x\) and \(y\) be any two points of \(A \cup B\) , and let \(Z \in A \cap B\) . Then \((x, z) \in V\) and \((z, y) \in V\) , so that \((x, y) \in \mathring{V}\) .

DEFINITION 2. A filter \(\mathfrak{F}\) on a uniform space \(\mathbf{X}\) is a Cauchy filter if for each entourage V of X there is a subset of X which is V-small and belongs to \(\mathfrak{F}\) .

Here again we may make our language more expressive by the use of the expressions ``sufficiently small set'' and ``a set as small as we please''; thus Definition 2 can be restated by saying that a Cauchy filter is one containing arbitrarily small sets.

An infinite sequence \((u_n)\) of points of a uniform space \(X\) is said to be a Cauchy sequence if the elementary filter associated with the sequence is a Cauchy filter. It comes to the same thing to say that for each entourage \(V\) of \(X\) there is an integer \(n_0\) such that for all integers \(m \geqslant n_0\) and \(n \geqslant n_0\) we have \((u_m, u_n) \in V\) .

PROPOSITION 2. On a uniform space X every convergent filter is a Cauchy filter.

If \(x\) is any point of \(X\) and \(V\) is any symmetric entourage of \(X\) , then the neighbourhood \(V(x)\) of \(x\) is \(V\) -small. If \(\mathfrak{F}\) is a filter which converges to \(x\) , there is a set of \(\mathfrak{F}\) contained in \(V(x)\) , and therefore \(V\) -small.

Clearly every filter which is finer than a Cauchy filter is a Cauchy filter.

PROPOSITION 3. Let \(f: \mathbf{X} \to \mathbf{X}'\) be a uniformly continuous mapping. Then the image under \(f\) of a Cauchy filter base on \(\mathbf{X}\) is a Cauchy filter base on \(\mathbf{X}'\) .

Let \(g = f \times f\) . If \(V'\) is an entourage of \(X'\) , then \(\overline{g}^1(V')\) is an entourage of \(X\) , and the image under \(f\) of a \(\overline{g}^1(V')\) -small set is \(V'\) -small; hence the result.

It follows in particular that if the uniformity of a uniform space X is replaced by a coarser uniformity, then every Cauchy filter with respect to the original uniformity remains a Cauchy filter with respect to the new uniformity.

This fact can be easily remembered in the following form: the finer the uniformity, the fewer Cauchy filters there are.

PROPOSITION 4. Let X be a set, let \((\mathbf{Y}_i)_{i\in I}\) be a family of uniform spaces, and for each \(i\in I\) let \(f_i\) be a mapping of X into \(\mathbf{Y}_i\) . Let X carry the coarsest uniformity \(\mathcal{U}\) for which the \(f_i\) are uniformly continuous. Then in order that a filter base \(\mathcal{B}\) on X should be a Cauchy filter base it is necessary and sufficient that \(f_i(\mathcal{B})\) should be a Cauchy filter base on \(\mathbf{Y}_i\) , for each \(i\in I\) .

The condition is necessary by Proposition 3. Conversely, suppose that it is satisfied, and let \(\mathbf{U}(\mathbf{V}_{\iota_1},\ldots ,\mathbf{V}_{\iota_n})\) be an entourage of the uniformity \(\mathfrak{U}\) {[}§ 2, no. 3, formula (1){]}. By hypothesis, for each index \(k\) there is a set \(\mathbf{M}_k\in \mathfrak{B}\) such that \(f_{\iota_k}(\mathbf{M}_k)\) is \(\mathbf{V}_{\iota_k}\) -small ( \(\iota \leqslant k\leqslant n\) ). Let \(\mathbf{M}\) be a set of \(\mathfrak{B}\) contained in \(\mathbf{M}_k\) for \(\iota \leqslant k\leqslant n\) ; then for each pair of points \(x,x^{\prime}\) of \(\mathbf{M}\) we have \([f_{\iota_k}(x),f_{\iota_k}(x')]\in \mathbf{V}_{\iota_k}\) for \(\iota \leqslant k\leqslant n\) , so that

\[
(x, x ^ {\prime}) \in \mathrm{U} (\mathrm{V} _ {t _ {1}}, \dots , \mathrm{V} _ {t _ {n}}).
\]

This completes the proof.

COROLLARY 1. If a Cauchy filter on a uniform space X induces a filter on a subset A of X, then this filter is a Cauchy filter on the uniform subspace A.

COROLLARY 2. A filter base \(\mathfrak{B}\) on a product \(\prod_{i\in I}X_i\) of uniform spaces is a Cauchy filter base if and only if, for each \(i\in I\) , \(\operatorname{pr}_i(\mathfrak{B})\) is a Cauchy filter base on \(X_i\) .

